\documentclass[a4paper,11pt]{article}
\usepackage[T1]{fontenc}
\usepackage[utf8]{inputenc}
\usepackage{lmodern,microtype}
\usepackage[margin=24mm,headheight=14pt]{geometry}
\usepackage{amsmath,amssymb,bm,mathtools}
\usepackage{booktabs,longtable,array,tabularx}
\usepackage{enumitem,xcolor,fancyhdr,listings}
\usepackage{hyperref}
\definecolor{ink}{HTML}{17364A}
\definecolor{accent}{HTML}{176C80}
\hypersetup{colorlinks=true,linkcolor=ink,citecolor=accent,urlcolor=accent,pdftitle={GW-QSense: Audited Research and Implementation Plan}}
\pagestyle{fancy}\fancyhf{}\fancyhead[L]{\small GW-QSense: audited plan}\fancyhead[R]{\small 18 September 2026}\fancyfoot[C]{\thepage}
\setlength{\parindent}{0pt}\setlength{\parskip}{5pt}
\setlist{itemsep=3pt,topsep=4pt,leftmargin=*}
\setlength{\emergencystretch}{2em}
\widowpenalty=10000
\clubpenalty=10000
\displaywidowpenalty=10000
\newcolumntype{P}[1]{>{\raggedright\arraybackslash}p{#1}}
\newcommand{\dd}{\mathrm{d}}
\newcommand{\Tr}{\operatorname{Tr}}
\newcommand{\Var}{\operatorname{Var}}
\newcommand{\ac}[2]{\{#1,#2\}}
\newcommand{\ket}[1]{\lvert #1\rangle}
\newcommand{\bra}[1]{\langle #1\rvert}
\newcommand{\hp}{h_{\mathrm P}}
\newcommand{\calE}{\mathcal E}
\newcommand{\eps}{\varepsilon}
\newcommand{\code}[1]{\texttt{#1}}
\lstset{basicstyle=\ttfamily\small,breaklines=true,columns=fullflexible,frame=single,showstringspaces=false}
\title{\color{ink}\bfseries GW-QSense\\[5pt]\Large Audited Research and Implementation Plan\\[6pt]\normalsize Gravitational-wave response, elastic transduction, and NV readout}
\author{Prepared for Daniel Junior Lifenya Fondo\\\small Critical synthesis of the four supplied inputs}
\date{18 September 2026}
\begin{document}
\maketitle

\begin{abstract}
Build a validated gravitational-wave-driven mechanical-response simulator first. Add NV spin readout only through an explicitly calculated local crystal strain and a documented material Hamiltonian. Keep the Dirac/Foldy--Wouthuysen (FW) investigation as a separate theoretical validation branch. The supplied plans contain correct ingredients, but also invalid operator projections, an incorrect Ward-identity proof, forbidden atomic transitions, inconsistent normalizations, and unsupported sensitivity claims. This document replaces those claims with derivations, stated approximation regimes, source corrections, and a staged implementation sequence. It establishes a defensible simulation programme; it does not establish a feasible NV gravitational-wave detector or certify an unseen complete FW calculation.
\end{abstract}

\textbf{Decision.} Retain the mechanical-to-NV programme as a conditional feasibility study. Reject the assertion that the pasted spin-1 manuscript has already derived its final NV gravitational Hamiltonian. Do not use its abstract or conclusion as statements of demonstrated results.

\textbf{Start here.} Implement the single-mode equation in Eq.~\eqref{eq:mode}, reproduce its analytic response and finite-duration limits, then calculate a noise budget. A functioning, documented V1 is the first finished product. Hydrogen, full FW calculations, hybrid phonon dynamics, and browser visualization are gated additions.

\textbf{Audit boundary.} Every distinct displayed calculation and substantive proposed extension in the four inputs is addressed below; repeated passages are treated once. Some claims can be disproved, some independently reproduced, and some only classified as unestablished because essential derivations or device specifications are absent. ``Checked'' does not mean ``all claims verified.'' The supplied references were checked for identity and for support of the claims for which they were used. The complete thesis, complete nine-term FW derivation, and an experimental device design are not part of these four inputs and are not certified by this audit.

\clearpage
\begingroup
\small
\setlength{\parskip}{0pt}
\tableofcontents
\endgroup
\clearpage

\section{One project, with explicit scientific boundaries}
\label{sec:scope}

\textbf{Question:} Given incident GW polarizations, a specified mechanical structure, and a specified readout, what signal and uncertainty result during a finite observation?

The common object is the \emph{detector response}, not an assumed universal spin Hamiltonian. The leading nonrelativistic mechanical response follows directly from tidal acceleration. A Dirac calculation is useful for investigating relativistic and spin corrections, but is not a prerequisite for deriving the bar response.

\begin{tabularx}{\textwidth}{@{}P{28mm}XX@{}}
\toprule
Layer & Input and output & Scientific status \\
\midrule
GW geometry & $\gamma_{+},\gamma_\times$, orientation $\to$ tidal tensor & Established linearized GR, within a stated frame and wavelength limit.\\
Elastic response & Tidal forcing $\to q_n,\dot q_n,\eps_{ij}$ & Established continuum mechanics; actual geometry and losses remain device inputs.\\
NV response & Local crystal strain $\to$ spin Hamiltonian and readout & Symmetry-based effective model; coefficients require provenance and a basis convention.\\
Inference & Signal plus noise $\to$ likelihood, SNR, uncertainty & Conditional on the stated measurement, resources, and noise.\\
Dirac/FW branch & Covariant matter model $\to$ controlled low-energy operators & Research audit; no unverified residual spin term enters detector claims.\\
\bottomrule
\end{tabularx}

An NV embedded in diamond does not automatically sense the strain of an external metal bar. A bonded diamond, support, interface, or strain concentrator needs its own mechanical transfer function, losses, and thermal noise. For an all-diamond device, use its actual anisotropic elastic response and feasible geometry.

\textbf{Success criteria.} Numerical accuracy, reproducibility, and a defensible feasibility result are separate outcomes. A negative feasibility result can be a useful calculation, but neither publication novelty nor career value follows automatically from the number of modules or visualizations.

\section{Conventions that must precede implementation}
\label{sec:conv}

Use SI units in the detector code. Adopt $x^0=ct$, signature $(-,+,+,+)$, and
\begin{equation}
 R^\rho{}_{\sigma\mu\nu}=\partial_\mu\Gamma^\rho_{\nu\sigma}-\partial_\nu\Gamma^\rho_{\mu\sigma}
 +\Gamma^\rho_{\mu\lambda}\Gamma^\lambda_{\nu\sigma}
 -\Gamma^\rho_{\nu\lambda}\Gamma^\lambda_{\mu\sigma}.
\end{equation}
Define the physical dimensionless perturbation $\gamma_{\mu\nu}=g_{\mu\nu}-\eta_{\mu\nu}$. If a canonical graviton field is introduced, write $\gamma_{\mu\nu}=\kappa h^{\rm can}_{\mu\nu}$ explicitly. In a common natural-unit convention $\kappa=\sqrt{32\pi G}$; do not copy that expression into SI code without restoring its dimensions.

The SI tidal tensor is
\begin{equation}
 \boxed{\calE_{ij}\equiv c^2R_{0i0j}=-\frac12\ddot\gamma^{\rm TT}_{ij}},
 \qquad \ddot X^i=-\calE^i{}_jX^j .
 \label{eq:tidal}
\end{equation}
$R_{0i0j}$ has units $\mathrm{m}^{-2}$, whereas $\calE_{ij}$ has units $\mathrm{s}^{-2}$. The inputs' $R_{0i0j}=-\ddot h_{ij}/2$ is correct only with $c=1$, a time-index convention that absorbs $c$, or an explicit redefinition as $\calE$.

For orthonormal transverse vectors $\bm e_1,\bm e_2$, set
\begin{equation}
 e^+_{ij}=e_{1i}e_{1j}-e_{2i}e_{2j},\quad
 e^\times_{ij}=e_{1i}e_{2j}+e_{2i}e_{1j},\quad
 \gamma^{\rm TT}_{ij}=\gamma_+e^+_{ij}+\gamma_\times e^\times_{ij}.
\end{equation}
Then $e^A_{ij}e^{Bij}=2\delta^{AB}$, both tensors are trace-free and transverse, and rotation of the polarization basis mixes plus and cross through twice the rotation angle. For a bar along $\bm n$, the driving strain is
\begin{equation}
 \gamma_{nn}=n^in^j\gamma^{\rm TT}_{ij}=F_+\gamma_++F_\times\gamma_\times,
 \qquad F_A=n^in^je^A_{ij}.
\end{equation}
There is no additional Michelson-detector factor $1/2$ in this bar definition.

\begin{tabularx}{\textwidth}{@{}P{31mm}XX@{}}
\toprule
Symbol & Meaning & Units/convention\\
\midrule
$\gamma_{ij}$ & Metric perturbation & Dimensionless; distinct from material strain.\\
$\eps_{ij}$ & Local crystal strain & Symmetric tensor, dimensionless. Tensor shear $\eps_{xy}$ is half engineering shear.\\
$\hp=2\pi\hbar$ & Planck constant & $\mathrm{J\,s}$; never use bare $h$ for three meanings.\\
$S_i$ & Spin-1 matrices & Dimensionless; $[S_i,S_j]=i\epsilon_{ijk}S_k$, $S^2=2I$.\\
$\nu_D,\gamma_e$ & NV splitting, Zeeman slope & Hz and Hz/T: $H/\hp=\nu_DS_z^2+\gamma_e\bm B\cdot\bm S+\cdots$.\\
$\Gamma_m, Q$ & Mechanical energy damping, quality factor & $\Gamma_m=\omega_0/Q$; amplitude decays at $\Gamma_m/2$.\\
\bottomrule
\end{tabularx}

Thus a Hamiltonian stored in Hz obeys $\dot\psi=-i2\pi(H/\hp)\psi$. A coefficient in Hz/strain must not be divided by $\hbar$ as though it were an energy.

\section{The corrected mechanical model}
\label{sec:mechanics}

\subsection{Derivation with a fixed mode normalization}

Assume a uniform slender bar, length $L$, area $A$, density $\rho$, total mass $M=\rho AL$, and effective longitudinal Young modulus $Y$. Its ends are free: $\partial_x\xi(\pm L/2,t)=0$. The reference origin is its centre. The leading GW force density in the long-wavelength, slow-motion approximation is
\begin{equation}
 \rho A\,\partial_t^2\xi-YA\,\partial_x^2\xi
 =\frac12\rho A\,x\,\ddot\gamma_{nn}(t).
 \label{eq:pde}
\end{equation}
Gravity multiplying the already GW-induced displacement is second order in the drive and is omitted here. Supports and other accelerations require additional forces and boundary conditions.

Choose the first nonzero longitudinal mode
\begin{equation}
 \xi(x,t)=q(t)u(x),\qquad u(x)=\sin(\pi x/L),\qquad
 \omega_0=\frac{\pi}{L}\sqrt{Y/\rho}.
\end{equation}
Here $q$ is an end-displacement amplitude, not a mass-normalized coordinate. Direct integration gives
\begin{align}
 m_1&=\rho A\int_{-L/2}^{L/2}u^2\dd x=\frac M2,\label{eq:mass}\\
 J&=\rho A\int_{-L/2}^{L/2}xu\dd x=\frac{2ML}{\pi^2},\label{eq:J}\\
 k_1&=YA\int_{-L/2}^{L/2}(u')^2\dd x=m_1\omega_0^2.
\end{align}
The bare overlap $\int x\sin(\pi x/L)\dd x=2L^2/\pi^2$ in the plans is correct. Its interpretation requires the mass and mode normalization above.

The tidal potential is $\frac12\int\rho A\calE_{xx}(x+\xi)^2\dd x$. Its term linear in $q$ is
\begin{equation}
 H_{\rm GW}=-\frac J2\ddot\gamma_{nn}(t)q=-\frac{ML}{\pi^2}\ddot\gamma_{nn}(t)q.
 \label{eq:Hbar}
\end{equation}
Consequently,
\begin{equation}
 \boxed{\ddot q+\Gamma_m\dot q+\omega_0^2q
 =\Lambda\ddot\gamma_{nn}(t),\qquad
 \Lambda\equiv\frac{J}{2m_1}=\frac{2L}{\pi^2}.}
 \label{eq:mode}
\end{equation}
This is the equation to implement first. The source convention for the sign of curvature or the definition of $q$ may reverse the displayed interaction sign; the acceleration, force, and Hamiltonian must remain mutually consistent. The magnitude agrees with the bar coupling of Tobar et al.~\cite{tobar} after matching $m_1=M/2$.

\subsection{Response, strain, and energy checks}

Use inverse Fourier components $e^{-i\omega t}$. The steady-state response is
\begin{equation}
 \mathcal H_{q\gamma}(\omega)=\frac{q(\omega)}{\gamma_{nn}(\omega)}
 =-\frac{\Lambda\omega^2}{\omega_0^2-\omega^2-i\Gamma_m\omega}.
 \label{eq:transfer}
\end{equation}
The strain in this one-dimensional model is
\begin{equation}
 \eps_{xx}(x,t)=q(t)\frac\pi L\cos(\pi x/L).
 \label{eq:localstrain}
\end{equation}
It vanishes at the free ends and is maximal at the centre. An NV at an end-displacement antinode does not thereby sit at a strain antinode. Real diamond needs transverse strains, orientation, and boundary conditions, not just this scalar field.

At the centre and at resonance,
\begin{equation}
 |\mathcal H_{\eps\gamma}(\omega_0)|=\frac{2Q}{\pi},\quad
 |q|=\Lambda Q\gamma_0,
 \quad \Delta f_{\rm power}=\frac{f_0}{Q},\quad
 \tau_{\rm amplitude}=\frac{2Q}{\omega_0}.
 \label{eq:resonance}
\end{equation}
The linewidth expression is the high-$Q$ power full width at half maximum. A one-mode model is useful near its isolated resonance and as a controlled truncation; do not extrapolate its high-frequency limit across higher resonances.

Track the mechanical energy and its balance:
\begin{equation}
 E_m=\frac12m_1\dot q^2+\frac12m_1\omega_0^2q^2,\qquad
 \dot E_m=\frac J2\ddot\gamma_{nn}\dot q-m_1\Gamma_m\dot q^2.
 \label{eq:energy}
\end{equation}
Energy should be conserved only when drive and damping vanish. A driven run should satisfy this balance, not an energy-conservation assertion.

\subsection{Finite signals: the essential feasibility correction}

For an initially resting undamped mode driven on resonance with
$\gamma_{nn}=\gamma_0\cos\omega_0t$ for $t\geq0$,
\begin{equation}
 q(t)=-\frac{\Lambda\omega_0\gamma_0}{2}t\sin\omega_0t,
 \qquad |\eps_{xx}(0,t)|_{\rm envelope}=\frac{\omega_0t}{\pi}\gamma_0.
 \label{eq:short}
\end{equation}
For high $Q$, the slowly varying resonant strain envelope is approximately
\begin{equation}
 \frac{2Q}{\pi}\gamma_0\bigl(1-e^{-\Gamma_mt/2}\bigr).
\end{equation}
This reproduces Eq.~\eqref{eq:short} for $t\ll2/\Gamma_m$. A chirp requires convolution or time integration with its actual phase; duration alone is not a substitute for residence near resonance. Ringdown after the signal can still be measured and must be included in an appropriate likelihood.

For arbitrary forcing, a useful independent reference is
\begin{equation}
 q(t)=q_{\rm hom}(t)+\Lambda\int_0^t
 \frac{e^{-\Gamma_m(t-s)/2}\sin[\omega_d(t-s)]}{\omega_d}
 \ddot\gamma_{nn}(s)\dd s,\quad
 \omega_d^2=\omega_0^2-\Gamma_m^2/4.
\end{equation}
This expression assumes an underdamped mode; use the corresponding overdamped or matrix-exponential solution outside that regime.

\textbf{Numerical illustration, not a proposed device.}
Let $M=1800\,\mathrm{kg}$, $v_s=5000\,\mathrm{m/s}$, $f_0=150\,\mathrm{Hz}$, and $Q=10^{10}$. A free longitudinal bar then has $L=v_s/(2f_0)=16.67\,\mathrm m$, not an unspecified independently chosen length. Its amplitude ring-up time is $2.122\times10^7\,\mathrm s=245.6$ days and its linewidth is $1.5\times10^{-8}\,\mathrm{Hz}$. For a $0.2$ s resonant coherent drive, the early strain gain is only $2f_0t=60$, versus the steady-state $6.37\times10^9$. These numbers expose why arbitrary $Q$ enhancement is not a burst sensitivity calculation.

\section{GW strain is not crystal strain}
\label{sec:strain}

The relevant material quantity is
\begin{equation}
 \eps^{\rm lab}_{ij}=\frac12(\partial_i u_j+\partial_j u_i),\qquad
 \eps^{\rm NV}=R\eps^{\rm lab}R^T,
 \label{eq:rotate}
\end{equation}
where the rows of $R$ are the NV basis vectors expressed in laboratory components. In detector coordinates this is the leading local elastic strain; curvature corrections to local spatial distances are beyond the retained small-device order.

For ideal freely falling test masses with appropriate initial conditions, the leading fractional proper-length change is $\delta L/L=\gamma_{nn}/2$. A bound elastic solid is not that freely falling configuration. Its response is frequency-dependent and constrained by stiffness and boundaries. Therefore neither $\eps_{ij}=\gamma_{ij}$ nor $\eps_{ij}=\gamma_{ij}/2$ is a universal constitutive law for diamond.

For multiple modes and an interface, document an actual transfer map
\begin{equation}
 \eps^{\rm NV}_{ij}(\omega)
 =\sum_A\mathcal T_{ij,A}(\omega)\gamma_A(\omega)
 +\eps^{\rm noise}_{ij}(\omega).
\end{equation}
This map includes mechanical mode shapes, GW overlaps, losses, NV position and orientation, and any bonded-material transfer. A metal-bar response cannot be combined with a diamond NV coupling constant without this link.

A compact longitudinal diamond device illustrates the frequency mismatch: the stipulated values $v_s=1.75\times10^4\,\mathrm{m/s}$ and $L=1\,\mathrm{mm}$ give $f_0=8.75\,\mathrm{MHz}$. Flexural modes can be much lower, but then the mode shapes, tidal overlap, and strain field must be rederived. Do not use the free longitudinal sine mode for a cantilever.

\section{The NV model that is justified}
\label{sec:nv}

\subsection{Spin matrices and the immediate correction}

In the order $\{\ket{+1},\ket0,\ket{-1}\}$,
\begin{equation}
 S_x=\frac1{\sqrt2}\begin{pmatrix}0&1&0\\1&0&1\\0&1&0\end{pmatrix},\quad
 S_y=\frac1{\sqrt2}\begin{pmatrix}0&-i&0\\i&0&-i\\0&i&0\end{pmatrix},\quad
 S_z=\begin{pmatrix}1&0&0\\0&0&0\\0&0&-1\end{pmatrix}.
\end{equation}
Therefore
\begin{equation}
 Q_+=S_x^2-S_y^2=\begin{pmatrix}0&0&1\\0&0&0\\1&0&0\end{pmatrix},\qquad
 Q_\times=\ac{S_x}{S_y}=\begin{pmatrix}0&0&-i\\0&0&0\\i&0&0\end{pmatrix}.
 \label{eq:Q}
\end{equation}
The restriction of $Q_+$ to $\{\ket0,\ket{+1}\}$ is the zero matrix, not $\mathrm{diag}(1,0)$. It couples $\ket{+1}$ and $\ket{-1}$, and annihilates $\ket0$. A simulation initialized in $\ket0$ with only $S_z^2$, $S_z$, and $Q_+$ will show no population response to that strain drive. Preparing a suitable state is part of the sensing protocol.

\subsection{A complete linear strain model}

Write $K=H/\hp$ in Hz. A symmetry-consistent $C_{3v}$ ground-triplet model is
\begin{equation}
 K=\nu_DS_z^2+\gamma_e\bm B\cdot\bm S+K_{\eps0}+K_{\eps1}+K_{\eps2}+K_{\rm control}.
 \label{eq:nv}
\end{equation}
Here $\nu_D\simeq2.87\times10^9\,\mathrm{Hz}$ and the positive Zeeman slope is $\gamma_e\simeq28\times10^9\,\mathrm{Hz/T}$. In the axes and tensor-strain conventions of Udvarhelyi et al.~\cite{udvarhelyi}, with their coupling symbols renamed $c_{ab}$,
\begin{align}
 K_{\eps0}&=[c_{41}(\eps_{xx}+\eps_{yy})+c_{43}\eps_{zz}]S_z^2,\\
 K_{\eps1}&=\tfrac12[c_{26}\eps_{zx}-\tfrac12c_{25}(\eps_{xx}-\eps_{yy})]\ac{S_x}{S_z}
 +\tfrac12(c_{26}\eps_{yz}+c_{25}\eps_{xy})\ac{S_y}{S_z},\\
 K_{\eps2}&=\tfrac12[c_{16}\eps_{zx}-\tfrac12c_{15}(\eps_{xx}-\eps_{yy})](S_y^2-S_x^2)
 +\tfrac12(c_{16}\eps_{yz}+c_{15}\eps_{xy})\ac{S_x}{S_y}.
 \label{eq:fullstrain}
\end{align}
All six $c_{ab}$ have units Hz/strain. The three parts connect $\Delta m_s=0$, $\pm1$, and $\pm2$, respectively in the unmixed $S_z$ basis. Thus $\Delta m_s=\pm1$ does not uniquely identify a vector gravitational interaction. Static state mixing further changes selection-rule language in the energy eigenbasis.

The DFT values reported in that source are, in MHz/strain,
\begin{center}
\begin{tabular}{@{}rrrrrr@{}}\toprule
$c_{41}$&$c_{43}$&$c_{25}$&$c_{26}$&$c_{15}$&$c_{16}$\\\midrule
$-6420\pm90$&$2300\pm200$&$-2600\pm80$&$-2830\pm70$&$5700\pm200$&$19660\pm90$\\\bottomrule
\end{tabular}
\end{center}
These are calculated parameters with the source's uncertainties, not six independently measured constants. Multiplying by $10^6$ converts this table to the code's Hz convention. Axis changes and engineering-shear conventions must transform the Hamiltonian, not just relabel parameters.

The measured reduced-model susceptibilities $d_\parallel=13.3\pm1.1$ and $d_\perp=21.5\pm1.2$ GHz/strain from Ovartchaiyapong et al.~\cite{ovart} refer to their geometry and effective strain convention. They must not be inserted as arbitrary members of the six-parameter tensor. The same work estimates a local axial strain sensitivity of about $3\times10^{-6}/\sqrt{\mathrm{Hz}}$ for that experiment. This is a dated, device-specific benchmark, not a universal present-day limit or an incident-GW sensitivity.

\subsection{Choose a sensing mechanism before choosing pulses}

For phase sensing, start with an axial strain term on a prepared $\ket0$--$\ket{+1}$ coherence, followed by a specified analysis pulse and fluorescence model. A longitudinal strain modulation changes relative phase but does not itself drive Rabi oscillations between those two states. With an ideal toggling function $y(t)=\pm1$,
\begin{equation}
 \phi(\theta)=2\pi\int_0^\tau y(t)\,\delta\nu(t;\theta)\dd t.
 \label{eq:phase}
\end{equation}
Hahn echo is the first pulse benchmark. Add CPMG or XY8 only after computing the filter for this actual Hamiltonian. For transverse or noncommuting couplings use the transformed operator $U_c^\dagger H_{\rm sig}U_c$, not an unjustified scalar toggling sign. Finite pulse durations, pulse errors, and signal suppression by the control must be included when relevant.

For resonant double-quantum control, define
\begin{equation}
 K_{\rm DQ}=\frac{\nu_s}{2}\sigma_z+d_{\rm eff}\eps_0\cos(2\pi f t)\sigma_x
\end{equation}
in the $\{\ket{+1},\ket{-1}\}$ subspace. At $f=\nu_s$, for weak drive and the rotating-wave approximation,
\begin{equation}
 P_{-1\to+1}(t)=\sin^2(\pi d_{\rm eff}\eps_0t),\quad
 f_R=|d_{\rm eff}\eps_0|,\quad \Omega_R=2\pi f_R.
 \label{eq:rabi}
\end{equation}
With negligible static transverse splitting, $\nu_s=2|\gamma_eB_z|$. Matching $150$ Hz would require $B_z=2.68$ nT in this idealization. Static strain, hyperfine couplings, and magnetic noise cannot be assumed negligible at such a small splitting. Microwave dressing or phase sensing can address a different frequency regime, but require a new explicit model.

The arithmetic $10^{10}\times10^{-21}=10^{-11}$ is correct. Under the \emph{stipulated local} strain $\eps_0=10^{-21}$ it gives $f_R=10^{-11}$ Hz, $\Omega_R=6.28\times10^{-11}\,\mathrm{rad/s}$, and a $\pi$-pulse time of about $1584$ years. It does not predict local strain from a GW. $N$ independent NV centres in the same classical drive retain this single-spin Rabi frequency; their ideal independent measurement information can scale as $N$. A $\sqrt N$ collective coupling requires a specified common-mode quantum interaction and collective state, not an ensemble-size substitution into Eq.~\eqref{eq:rabi}.

\subsection{What a material derivation would actually require}

A microscopic derivation needs an electronic many-body Hamiltonian with crystal potential, electron interactions, spin-orbit and spin-spin terms, its strain dependence, and a projection onto the NV ground manifold. Perturbation theory can then generate terms such as
\begin{equation}
 H_{\rm eff}=PH_0P+PVP+PVQ(E_0-QH_0Q)^{-1}QVP+\cdots,
\end{equation}
including mixed perturbations when $V$ has multiple physical origins. This is where material spin-strain coefficients arise. A rank argument alone neither computes them nor shows that a spin-independent momentum operator becomes a spin quadrupole.

\section{Quantizing the mechanical mode is an optional next step}
\label{sec:quantum}

For the normalization in Eq.~\eqref{eq:mass},
\begin{equation}
 q_{\rm zpf}=\sqrt{\frac{\hbar}{2m_1\omega_0}}=\sqrt{\frac{\hbar}{M\omega_0}},\quad
 \hat q=q_{\rm zpf}(b+b^\dagger),\quad
 \eps_{\rm zpf}(x)=u'(x)q_{\rm zpf}.
 \label{eq:zpf}
\end{equation}
Using total mass $M$ inside $\sqrt{\hbar/(2M\omega_0)}$ with the same unnormalized sine mode is a factor-$\sqrt2$ error. A mode rescaled to mass $M$ is allowed only if $q$, strain, and overlap are rescaled consistently.

The GW drive becomes
\begin{equation}
 H_{\rm GW}(t)=-\frac L{\pi^2}\sqrt{\frac{M\hbar}{\omega_0}}
 \ddot\gamma_{nn}(t)(b+b^\dagger).
 \label{eq:quantbar}
\end{equation}
For a harmonic GW, its energy-amplitude magnitude is
$\mathcal E_{\rm drive}=(L/\pi^2)\sqrt{M\hbar/\omega_0}\,\omega_{\rm GW}^2\gamma_0$.
The corresponding angular-frequency amplitude is $\mathcal E_{\rm drive}/\hbar$.
The input called the energy amplitude $g_{\rm GW}$ without fixing its units and used one $\omega$ for two frequencies. They coincide only on resonance.

A classical signal is not ``converted into zero-point strain.'' Zero-point strain sets the quantum matrix element; the classical drive creates a coherent displacement on top of the oscillator's fluctuations. For $H=\hbar\omega_0b^\dagger b+\lambda(t)(b+b^\dagger)$, initially in its ground state,
\begin{equation}
 \alpha_I(t)=-\frac i\hbar\int_0^t\lambda(s)e^{i\omega_0s}\dd s,
 \qquad P_n=e^{-|\alpha|^2}\frac{|\alpha|^{2n}}{n!}.
\end{equation}
This is a useful exact benchmark for the lossless driven oscillator. A phonon jump under a classical drive does not alone prove that gravity is quantized; the distinction is analysed explicitly by Carney, Domcke and Rodd~\cite{carney}.

For an embedded NV, the linear quantum interaction is
\begin{equation}
 H_{\rm NV-ph}=\hp\sum_{ij}\frac{\partial K_\eps}{\partial\eps_{ij}}
 \eps_{{\rm zpf},ij}(\bm r_{\rm NV})(b+b^\dagger),
\end{equation}
with independent tensor components counted consistently. Do not include the same coherent strain twice, once in $K_\eps(t)$ and again through an explicitly driven oscillator.

The bar illustration of Section~\ref{sec:mechanics} has $q_{\rm zpf}\simeq7.9\times10^{-21}$ m and central $\eps_{\rm zpf}\simeq1.49\times10^{-21}$. A stipulated $10^{10}$ Hz/strain readout would couple at only about $1.49\times10^{-11}$ Hz per zero-point amplitude, before interface and orientation factors. These numbers are a scale check for that bar, not a universal bound on nanomechanical devices.

Before constructing a large spin-phonon Hilbert space, compare thermal occupation
\begin{equation}
 \bar n_{\rm th}=\bigl[e^{\hbar\omega_0/(k_BT)}-1\bigr]^{-1}
\end{equation}
with the intended cutoff and compare coupling rates to mechanical and spin decoherence. At $150$ Hz, $\hp f_0/k_B\simeq7.2$ nK; at $1$ mK, $\bar n_{\rm th}\simeq1.39\times10^5$. A $30$-state oscillator is not an equilibrium thermal model there. Ground-state preparation and subsequent bath heating are separate assumptions. With $Q=10^{10}$, the initial ground-state heating rate $\Gamma_m\bar n_{\rm th}\simeq0.0131\,\mathrm{s}^{-1}$ still needs to be included.

\section{Sensitivity and inference: required before a detector claim}
\label{sec:noise}

\subsection{Classical response and noise first}

For viscous damping in the classical thermal regime $k_BT\gg\hbar\omega$, the one-sided force PSD per Hz is
\begin{equation}
 S_F^{\rm th}=4k_BT m_1\Gamma_m,\quad
 \chi_m(\omega)=\frac1{m_1(\omega_0^2-\omega^2-i\Gamma_m\omega)},\quad
 S_q^{\rm th}=|\chi_m|^2S_F^{\rm th}.
\end{equation}
This follows from fluctuation--dissipation for the stipulated damping law; structural damping requires a different frequency dependence. Let $S_q^{\rm imp}$ denote displacement-equivalent readout noise. If these contributions are independent,
\begin{equation}
 S_\gamma(\omega)=\frac{|\chi_m|^2S_F^{\rm th}+S_q^{\rm imp}+S_q^{\rm other}}
 {|\mathcal H_{q\gamma}|^2}.
 \label{eq:noise}
\end{equation}
Include correlated terms when noise channels are correlated. The thermal contribution at resonance reduces to
\begin{equation}
 S_\gamma^{\rm th}(\omega_0)=\frac{4k_BT\Gamma_m}{m_1\Lambda^2\omega_0^4}
 =\frac{2\pi^4 k_BT}{ML^2Q\omega_0^3}.
\end{equation}
The illustrative bar at $1$ mK gives about $8.0\times10^{-25}/\sqrt{\mathrm{Hz}}$ for this thermal-only input-referred ASD. It excludes readout, support, environmental, and preparation noise and says nothing by itself about realizability.

If the $3\times10^{-6}/\sqrt{\mathrm{Hz}}$ historical NV local-strain benchmark were usable with unit geometric projection at steady-state resonance, its readout-only input equivalent would be
\begin{equation}
 \eta_\gamma^{\rm imp}=\frac{\eta_\eps}{|\mathcal H_{\eps\gamma}|}
 =\frac{3\times10^{-6}}{2Q/\pi}
 =4.71\times10^{-16}/\sqrt{\mathrm{Hz}}\quad(Q=10^{10}).
\end{equation}
This is not a sensitivity forecast: it assumes an unproved combination of mode, interface, readout, bandwidth, and long coherent observation. An ensemble's $N^{-1/2}$ improvement is conditional on independent noise, readout efficiency, usable NV orientations, unchanged coherence, and adequate optical resources. It cannot be freely compounded with all optimistic mechanical assumptions.

A strain amplitude $\gamma_0$ and an ASD $\sqrt{S_\gamma}$ have different units. For a known finite waveform $s(t)$ in stationary Gaussian output noise, use
\begin{equation}
 \mathrm{SNR}^2=4\int_0^\infty\frac{|\tilde s(f)|^2}{S_n(f)}\dd f,
 \label{eq:snr}
\end{equation}
with a consistent Fourier transform, observation window, transient response, and PSD convention. Do not compare $10^{-21}$ directly to a number per $\sqrt{\rm Hz}$ or import an astrophysical audio-band amplitude into a MHz/GHz proposal without a source model.

\subsection{A measurement model before Fisher information}

Specify state preparation, interrogation time $\tau$, control, contrast $C$, decoherence, readout distribution, reset time, number of probes, and number of repetitions. For an ideal binary phase measurement one useful model is
\begin{equation}
 p(\theta)=\tfrac12[1+C e^{-\chi(\tau)}\cos(\phi(\theta)+\phi_a)],\qquad
 F_C=\frac{[\partial_\theta p]^2}{p(1-p)}.
\end{equation}
Operate at a calibrated phase slope. Real fluorescence readout generally requires its photon-count distribution; the binary expression is a benchmark, not an automatic NV likelihood.

For $\nu$ independent repetitions of a model whose full-probe information is $F_C$, the local unbiased Cramer--Rao bound is $\Var\hat\theta\geq1/(\nu F_C)$. If $N$ independent identically prepared spins are already included in $F_C$, do not multiply by $N$ again. Account for total time $T_{\rm tot}=\nu(\tau+t_{\rm dead})$. A scalar bound assumes other parameters are known; uncertain mechanical gain, calibration, and phase require nuisance parameters and the inverse information matrix. Without calibration, strain amplitude and gain can be unidentifiable.

For a normalized pure state the supplied formula is correct:
\begin{equation}
 F_Q=4\left(\langle\partial_\theta\psi|\partial_\theta\psi\rangle
 -|\langle\psi|\partial_\theta\psi\rangle|^2\right).
\end{equation}
But $H(\theta)=H_0+\theta G$ does not in general imply $F_Q=4\tau^2\Var(G)/\hbar^2$. The latter needs a commuting, time-independent generator. In general set
\begin{equation}
 \mathcal G_\theta=iU_\theta^\dagger\partial_\theta U_\theta,
 \qquad F_Q=4\Var_{\psi_0}(\mathcal G_\theta).
\end{equation}
For a mixed state with eigenvalues $\lambda_a$,
\begin{equation}
 F_Q=2\sum_{a,b:\lambda_a+\lambda_b>0}
 \frac{|\langle a|\partial_\theta\rho|b\rangle|^2}{\lambda_a+\lambda_b}.
\end{equation}
QFI is an upper limit on accessible measurement information, not proof that the proposed fluorescence measurement attains it.~\cite{braunstein,degen}

\textbf{Feasibility gate.} Publish a detector sensitivity only after providing one internally consistent parameter set, local transfer function, finite waveform or signal statistics, all dominant noise contributions, readout likelihood, total resource count, and uncertainty propagation. If this gate fails, report response coefficients and conditional limits, not a demonstrated detection threshold.

\section{What fails in the proposed Dirac-to-NV manuscript}
\label{sec:projection}

\subsection{The vector projection: correct conclusion, wrong intermediate value}

In a pure two-spin-$1/2$ triplet, with dimensionless spins,
\begin{equation}
 \bm S=\bm s_1+\bm s_2,\quad
 \bm s_1\cdot\bm s_2=\tfrac14,\quad
 \bm s_1\cdot\bm S=\tfrac34+\tfrac14=1.
\end{equation}
Thus the projection coefficient is $1/[S(S+1)]=1/2$, giving
\begin{equation}
 P\bm s_1P=P\bm s_2P=\tfrac12\bm S.
\end{equation}
The manuscript's intermediate value $\langle\bm s_a\cdot\bm S\rangle=1/2$ is false; it would instead give a coefficient $1/4$ in its own formula.

For spin-independent coefficients acting on separate orbital coordinates, spin projection yields
\begin{equation}
 P(\bm v_1\cdot\bm s_1+\bm v_2\cdot\bm s_2)P
 =\tfrac12(\bm v_1+\bm v_2)\cdot\bm S
\end{equation}
within that spin-only model. Replacing this by $\bm v\cdot\bm S$ requires equal coefficients, or a properly performed orbital average. It does not permit replacing two momentum-dependent orbital operators by an unspecified common classical field. Identical-electron antisymmetry and the NV orbital state must also be respected.

\subsection{The quadrupole identity is valid; its proposed origin is not}

For the tensor written in the input,
\begin{equation}
 T_{ij}=s_{1i}s_{2j}+s_{1j}s_{2i}-\tfrac23(\bm s_1\cdot\bm s_2)\delta_{ij},
\end{equation}
projection onto the triplet does give
\begin{equation}
 PT_{ij}P=\tfrac12\ac{S_i}{S_j}-\tfrac23\delta_{ij}I.
\end{equation}
This follows by expanding $\ac{S_i}{S_j}$ and using the spin-$1/2$ anticommutators. The factor $1/2$ in the input's equivalent expression is correct.

However, $p_ip_j$ is an orbital operator and spin-independent. In an orbital singlet with no spin-dependent admixture,
\begin{equation}
 P_{\rm orb}\,p_ip_j\,P_{\rm orb}=C_{ij}P_{\rm orb}
\end{equation}
acts as the identity on spin. It does not become $T_{ij}$ merely because both objects have rank two in some space. The tensor identity proves the form of a \emph{spin tensor}, not its coefficient or its generation from orbital kinetic energy. Spin-orbit and spin-spin physics or an independently calibrated effective model are required.

There is also a separate factor-of-two error: with an Einstein sum over all $i,j$,
\begin{equation}
 \Xi\gamma_{ij}\ac{S_i}{S_j}
 =2\Xi\gamma_+(S_x^2-S_y^2)
\end{equation}
for $\gamma_{xx}=-\gamma_{yy}=\gamma_+$. The input's last expression loses the factor two unless $\Xi$ is explicitly redefined. A general NV coupling is a fourth-rank strain-to-spin tensor with crystal symmetry, not necessarily a single scalar $\Xi$.

\subsection{The ``gravitomagnetic field'' is not an established NV drive}

The expression $\nabla\times(\gamma\cdot\bm p)$ is momentum-dependent. In a spatially homogeneous long-wavelength field its derivative contribution vanishes for commuting free momenta; at finite wavelength it needs ordering and orbital matrix elements. Replacing it by a classical oscillatory field assumes precisely the material reduction that has not been performed. It is not a magnetic field measured in tesla.

Changing from TT coordinates to a detector frame does not break a gauge symmetry or create an additional independent interaction. In Fermi normal coordinates about a nonrotating geodesic,~\cite{manasse}
\begin{align}
 g_{00}&=-1-R_{0i0j}X^iX^j+\cdots,\\
 g_{0i}&=-\tfrac23R_{0jik}X^jX^k+\cdots,\\
 g_{ij}&=\delta_{ij}-\tfrac13R_{ikjl}X^kX^l+\cdots.
\end{align}
The connection vanishes on the reference worldline but curvature does not. Finite-size spin/curvature effects, rotating frames, and supported laboratories require the corresponding tetrad, apparatus, and controls. It is unjustified both to invent a universal surviving vector drive and to assert that all gravitational spin effects vanish.

The input also labels a shift-vector/orbital term ``gravitoelectric (Lense--Thirring).'' Lense--Thirring frame dragging belongs to the gravitomagnetic sector; it is not synonymous with every nonzero $g_{0i}$ in a wave-adapted coordinate system. A weak static metric gives the rest-energy correction
\begin{equation}
 \delta H_{\rm rest}=-\tfrac12mc^2\gamma_{00}
\end{equation}
in the stated convention, not a universal term proportional to the trace $\eta^{\mu\nu}\gamma_{\mu\nu}$. A spatially varying tidal potential is not a removable common phase for orbital transitions. The input's trace mass term therefore cannot be accepted merely by labelling it irrelevant.

\section{Correct gauge and Hamiltonian equivalence checks}
\label{sec:ward}

\subsection{The Ward identity must remain four-dimensional}

Up to external-state normalization, a linear graviton vertex can be written
\begin{equation}
 \mathcal M=-\frac\kappa2\epsilon_{\mu\nu}\langle f|T^{\mu\nu}(q)|i\rangle.
\end{equation}
A full polarization shift $\delta\epsilon_{\mu\nu}=q_\mu\xi_\nu+q_\nu\xi_\mu$ then gives
\begin{equation}
 \delta\mathcal M=-\kappa\xi_\nu q_\mu\langle T^{\mu\nu}\rangle=0
\end{equation}
when the \emph{complete} matter stress tensor is conserved and the required boundary and on-shell conditions hold. But with $q^\mu=(\omega/c,\bm k)$,
\begin{equation}
 k_iT^{ij}=\frac\omega c T^{0j},
\end{equation}
not zero in general. Nonrelativistic motion does not justify deleting the energy/momentum-density term for a dynamical transition. A spatial-only shift also need not preserve the TT conditions. The repeated proof in the attachments is invalid.

A bound electron by itself exchanges momentum with its binding field and nucleus. Its kinetic stress is not the full conserved stress tensor. Include binding interactions, electromagnetic stresses, the nucleus and, where applicable, supports or momentum exchange with apparatus. A free massive particle cannot absorb one real graviton while remaining an on-shell free particle: the on-shell conditions require $p\cdot q=0$, incompatible with nonzero positive-frequency null $q$ and future timelike $p$.

The proposed formula $T^{ij}_{\rm spin}=\epsilon^{ikl}q_kS_l+O(q^2)$ is malformed: its right-hand side has no free $j$ index and is not a symmetric stress tensor. Its antisymmetry argument proves nothing about the omitted complete amplitude. Delete it and derive any spin stress from the specified action.

\subsection{The momentum-position identity is false}

For $H_0=\bm p^2/(2m)+V(\bm x)$,
\begin{equation}
 [H_0,[H_0,x_ix_j]]
 =-\frac{2\hbar^2}{m^2}p_ip_j
 +\frac{\hbar^2}{m}(x_i\partial_jV+x_j\partial_iV).
\end{equation}
Consequently, for $\omega_{fi}=(E_f-E_i)/\hbar$,
\begin{equation}
 \boxed{\langle p_ip_j\rangle_{fi}
 =\frac m2\langle x_i\partial_jV+x_j\partial_iV\rangle_{fi}
 -\frac{m^2\omega_{fi}^2}{2}\langle x_ix_j\rangle_{fi}.}
 \label{eq:commutator}
\end{equation}
The input's $\langle p_ip_j\rangle=m^2\omega_{fi}^2\langle x_ix_j\rangle$ has the wrong sign and factor and omits the binding force. For a harmonic oscillator $0\to2$, $\langle2|p^2|0\rangle=-m^2\omega_0^2\langle2|x^2|0\rangle$ while $\omega_{20}=2\omega_0$; this supplies a simple counterexample.

The correct long-wavelength stress-quadrupole relation for a localized conserved system is
\begin{equation}
 I^{ij}=\frac1{c^2}\int T^{00}x^ix^j\dd^3x,\qquad
 \int T^{ij}\dd^3x=\frac12\ddot I^{ij}.
\end{equation}
It follows from two integrations by parts and conservation, with surface terms absent. Between stationary states it becomes $\int T^{ij}_{fi}\dd^3x=-\omega_{fi}^2 I^{ij}_{fi}/2$. Hence the TT interaction $-\gamma_{ij}\int T^{ij}/2$ gives the on-resonance matrix element $+\omega_{fi}^2\gamma_{ij}I^{ij}_{fi}/4$, consistent with
\begin{equation}
 H_{\rm tidal}=\tfrac12\calE_{ij}I^{ij}=-\tfrac14\ddot\gamma_{ij}I^{ij}.
 \label{eq:tidalH}
\end{equation}
This is an observable equivalence under its assumptions, not permission to identify incomplete Hamiltonians term by term. For a two-body atom, use the reduced mass and relative coordinate at leading nonrelativistic order.

\subsection{A transparent free-particle transformation}

For homogeneous $\gamma_{ij}(t)$, start from $H_{\rm TT}=p^2/(2m)-\gamma_{ij}p_ip_j/(2m)$. To first order in $\gamma$, the generating function
\begin{equation}
 F_2(\bm x,\bm P,t)=x_i(\delta_{ij}+\gamma_{ij}/2)P_j
\end{equation}
produces $X=(I+\gamma/2)x$ and $H'=P^2/(2m)+\dot\gamma_{ij}X_iP_j/2$. A second generating function
\begin{equation}
 F'_2(\bm X,\bm\Pi,t)=\bm X\cdot\bm\Pi-\tfrac m4\dot\gamma_{ij}X_iX_j
\end{equation}
cancels the mixed momentum term and gives $H''=\Pi^2/(2m)-m\ddot\gamma_{ij}X_iX_j/4$. Quantum mixed products must be symmetrized. For a binding potential, transform the potential and its source too; states, preparation, and measured operators also transform.

A gauge visualization should compare these controlled endpoints and transformed observables. An arbitrary slider interpolating Hamiltonians is not a physical gauge transformation unless the accompanying time-dependent canonical/unitary map is supplied.

\subsection{What is still required for the FW branch}

For $H=\beta m+\mathcal E+\mathcal O$ in $c=\hbar=1$ units, a useful \emph{partial} ledger contains
\begin{equation}
 \frac{\beta\mathcal O^2}{2m},\quad
 -\frac{\beta\mathcal O^4}{8m^3},\quad
 -\frac{[\mathcal O,[\mathcal O,\mathcal E]]}{8m^2},\quad
 -\frac{i[\mathcal O,\dot{\mathcal O}]}{8m^2}.
\end{equation}
It is not automatically the complete $1/m^3$ result for general time-dependent backgrounds. Specify counting in $p/m$, $\hbar\omega/(mc^2)$, field amplitude, and spatial gradients, especially when an even perturbation itself scales with $m$. Include the time derivative of every unitary transformation. Verify the conserved inner product and any wavefunction rescaling before testing flat-measure Hermiticity. Anticommutators help with ordering but do not establish gauge invariance, self-adjoint domains, or completeness.

Quach's correction~\cite{quach} concerns an electromagnetic gauge-variant term under $A_i\mapsto A_i+\partial_i\chi$. It is not the spatial-only gravitational Ward proof offered in the inputs. Its no-precession statement also uses a stated small-gradient limit; it cannot be promoted to a general theorem forbidding all spin response.

For a controlled benchmark, compare $E=\sqrt{m^2c^4+c^2p^2}$ with $mc^2+p^2/(2m)-p^4/(8m^3c^2)$; the next term is $p^6/(16m^5c^4)$. Compare states and observables in the same representation. A time-dependent $4\times4$ momentum block tests a spatially homogeneous approximation, not gradient terms of a propagating wave, which require spatial resolution or momentum coupling. Removing spatial derivatives does not in general remove $\dot{\mathcal O}$ terms.

\section{Hydrogen benchmarks: retain the corrected calculation}
\label{sec:hydrogen}

Hydrogen is a separate analytic regression problem; it need not block the mechanical V1. Use the nonrelativistic Coulomb problem with reduced mass $\mu$ and corresponding Bohr radius $a$. In the usual Condon--Shortley phase convention,
\begin{equation}
 R_{21}(r)=\frac1{2\sqrt6\,a^{3/2}}\frac ra e^{-r/(2a)},\quad
 \langle r^2\rangle_{2p}=30a^2,\quad \langle r^{-1}\rangle_{2p}=\frac1{4a}.
\end{equation}
The angular factor for $\bra{2,1,+1}(x^2-y^2)\ket{2,1,-1}$ is $-2/5$, so
\begin{equation}
 \bra{2,1,+1}(x^2-y^2)\ket{2,1,-1}=-12a^2.
\end{equation}
A rephased basis can change an off-diagonal sign, not its magnitude. The competing $3a^2/2$ value in the inputs is wrong. For $\gamma_+(t)=\gamma_0\cos\omega t$, the projected perturbation is
\begin{equation}
 H_{2p}(t)=\begin{pmatrix}0&0&-C(t)\\0&0&0\\-C(t)&0&0\end{pmatrix},\quad
 C(t)=3\mu a^2\omega^2\gamma_+(t),
\end{equation}
with instantaneous eigenvalues $0,\pm C(t)$. These are not the input's $0,\pm(3/2)m_ea_0^2\omega^2\gamma_+$.

Within an exactly degenerate subspace and for fixed linear polarization, $H(t)$ is a scalar function times a fixed matrix and commutes with itself at different times. Over complete sinusoidal cycles its integrated first-order perturbation vanishes. Instantaneous diagonalization therefore does not establish resonant absorption. Finite-duration phase evolution or transitions between physically split levels must be treated separately. For an atomic observation include relevant fine structure, hyperfine structure, fields, and linewidths.

The corrected Gaunt identity is
\begin{align}
 \int Y_{l'm'}^*Y_{2q}Y_{lm}\dd\Omega
 ={}&(-1)^{m'}\sqrt{\frac{(2l'+1)5(2l+1)}{4\pi}}
 \begin{pmatrix}l'&2&l\\0&0&0\end{pmatrix}
 \begin{pmatrix}l'&2&l\\-m'&q&m\end{pmatrix}.
\end{align}
The input omitted complex-conjugation phase and the minus sign on $m'$. For axes aligned with GW propagation, plus/cross quadrupoles have $q=\pm2$, with $|l-l'|\leq2\leq l+l'$ and even $l+l'$. Thus $\Delta l=0,\pm2$ is shorthand subject to triangle constraints; $0\to0$ and $1s\to2p$ are forbidden. A rotated axis mixes tensor components and changes the allowed displayed $\Delta m$ values.

\subsection{Absorption rate and the Planck-area statement}

Write $H_1(t)=V\cos\omega t$ with
$V=\mu\omega^2\gamma_0(x^2-y^2)/4$. The positive-frequency transition amplitude is $V_{fi}/2$, so
\begin{equation}
 \Gamma_{i\to f}=\frac{2\pi}{\hbar}\left|\frac{V_{fi}}2\right|^2
 \delta(E_f-E_i-\hbar\omega).
\end{equation}
Using the full cosine amplitude inside the Golden Rule would overestimate the rate by four. For an isolated coherently driven discrete transition, finite-time perturbation theory or Rabi dynamics is more appropriate than interpreting a delta function as a finite rate.

For one real plus polarization with our tensor normalization, the cycle-averaged flux is
\begin{equation}
 I_{\rm GW}=\frac{c^3}{32\pi G}\omega^2\gamma_0^2.
\end{equation}
That supplied flux is correct for this convention. Let $Q_{fi}=\bra f(x^2-y^2)\ket i$. The direction- and polarization-specific spectral cross section is then
\begin{equation}
 \sigma_{fi}(\omega)=\frac{\pi^2G\mu^2\omega^3}{c^3}|Q_{fi}|^2
 \delta(E_f-E_i-\hbar\omega)
 =\frac{\pi^2G\mu^2\omega^3}{\hbar c^3}|Q_{fi}|^2\delta(\omega-\omega_{fi}).
 \label{eq:xs}
\end{equation}
A physical linewidth or finite interrogation window replaces the delta distribution. Peak, integrated, and averaged cross sections are different quantities.

The allowed $1s\to3d,m=2$ matrix element is $Q_{fi}=81a^2/128$ in this phase convention; the claimed nonzero $1s\to2p$ quadrupole matrix element must be deleted. Boughn and Rothman~\cite{boughn}, Eqs.~(3.12)--(3.16), average over directions and unpolarized waves and define
\begin{equation}
 \bar\sigma\equiv\frac1{\omega_{fi}}\int\sigma(\omega)\dd\omega
 =\frac{3^4\pi^2}{5\,2^9}\frac{G\hbar}{c^3}
 =0.31228045\,\ell_{\rm Pl}^2
\end{equation}
for that specified transition in their hydrogen approximation. Their tensor matrix elements satisfy $D_{xx}=-D_{yy}$, $|D_{xy}|=|D_{xx}|$, $D_{xx}=81a^2/256$, and
$\langle|D_{ij}e_{ij}|^2\rangle=8|D_{xx}|^2/5$ in their polarization convention. Substituting $\omega_{fi}=4\hbar/(9\mu a^2)$ reproduces the coefficient. Atomic constants cancel in this particular normalized average; the result is neither a universal peak cross section nor a measure of practical detector efficiency. Numerically it is approximately $8.16\times10^{-71}\,\mathrm m^2$.

\section{Execution sequence and stopping rules}
\label{sec:roadmap}

\subsection{V1: one finished native C++ response simulator}

The following order replaces the incompatible roadmaps. Calendar estimates are planning assumptions; passing a gate matters more than reaching a week number. A first block of roughly four to six focused weeks is a budget for V1, not a promise that the entire research programme can be finished in twelve weeks.

\begin{longtable}{@{}P{17mm}P{62mm}P{74mm}@{}}
\toprule
Stage & Build and learn & Completion gate\\\midrule\endhead
0 & One conventions document; waveform, mode, run-configuration structs; SI units; a CMake executable and test target. & Every parameter has units and provenance. Distinguish metric strain, crystal strain, Hz, rad/s, total mass and modal mass.\\
1 & Analytic sine and smooth finite burst; one damped-mode solver; CSV output of $t,\gamma,q,\dot q,\eps,E_m$. & Zero forcing, free oscillation, damping and analytic harmonic response pass; error decreases at the expected order.\\
2 & Derive mode mass, stiffness, overlap and local strain; add polarization/orientation. & Reproduce Eqs.~\eqref{eq:mass}--\eqref{eq:localstrain}; changing normalization leaves physical displacement unchanged.\\
3 & Finite waveform convolution or ODE response; ringdown; controlled CSV import. & Analytic and numerical results agree; chirps do not receive arbitrary steady-state gain; energy balance closes.\\
4 & Transfer magnitude and phase; one-sided PSD and a thermal/readout noise model. & Parseval/PSD normalization and injected-signal recovery work; all SNRs identify bandwidth, duration and noise.\\
5 & Reproducible command, benchmark configuration, plots, validation report and source references. & A fresh native build reproduces documented results. Tag V1 only now.\\
\bottomrule
\end{longtable}

Learn the C++ feature needed by the current stage: basic structs and functions, ownership and \code{std::vector}, then small matrix operations. There is no reason to begin with an extensible class hierarchy for every theoretical noun.

\textbf{First concrete task.} Write the dimensionless version with $s=\omega_0t$, $y=q/L$:
\begin{equation}
 y''+Q^{-1}y'+y=\frac2{\pi^2}\gamma_{nn}''(s).
\end{equation}
For $\gamma_{nn}=\gamma_0\cos(rs)$, check the forcing
$-(2/\pi^2)r^2\gamma_0\cos(rs)$ against a complex analytic solution. Scale the state by $\gamma_0$ for tiny signals and convert back afterward; never let default absolute ODE tolerances erase a $10^{-21}$ response.

A small repository is enough initially:

\begin{tabularx}{\textwidth}{@{}P{52mm}X@{}}\toprule
Path/module & Responsibility\\\midrule
\code{include/gwq/waveform.hpp} & Values, analytic derivatives or a documented sampled-waveform policy.\\
\code{include/gwq/bar.hpp} & Mode parameters, response equation, local strain and energy.\\
\code{src/main.cpp} & Read configuration, run one calculation, export results.\\
\code{tests/} & Independent analytic comparisons and error convergence.\\
\code{python/reference.py} & Reference quadrature, convolution, and plots.\\
\code{configs/}, \code{docs/} & Reproducible inputs, assumptions, validation and sources.\\
\bottomrule\end{tabularx}

C++20, CMake, a test framework, and small linear-algebra support when needed are sufficient. Add JSON for configuration when it solves a real interface need. Defer HDF5, OpenMP, GPU work, generic FEM, Docker, and a frontend until data volume or measured runtime justifies them. CMake generates/invokes native build workflows; it is not itself the C++ compiler. Containers help pin environments but do not guarantee identical floating-point results across every platform. The finance/hedging appendix is unrelated to this project and is removed.

\subsection{V2: calibrated NV phase readout}

Add the spin-1 matrices, the complete or explicitly restricted strain model, an NV-axis transform, a static spectrum, state preparation, and one readout protocol. First reproduce the relevant observable from an established mechanical NV experiment~\cite{ovart,teissier,macquarrie,barfuss} in its own geometry. This is a test of the NV interface; it does not validate the proposed GW-to-device transfer.

Then supply the actual resonator-to-diamond strain and propagate its uncertainty. Start with classical strain driving the NV. Include static strain, thermal and magnetic noise, and hyperfine structure whenever they are not small compared with the relevant linewidth or splitting. An electron-only qutrit is an approximation, not the full NV system.

\textbf{Gate:} Can one documented geometry, measurement, and noise model support a useful incident-strain estimate? If not, finish the conditional feasibility report and stop detector optimization. Do not add more pulses or more NVs solely to make a target number appear.

\subsection{V3 and research extensions, only after the relevant gates}

\begin{longtable}{@{}P{39mm}P{46mm}P{68mm}@{}}\toprule
Idea from the inputs & Decision & Required condition\\\midrule\endhead
Hydrogen $2p$ and allowed absorption & Retain as compact analytic companion. & Use Section~\ref{sec:hydrogen}; do not put it on V1's critical path.\\
Full Dirac/FW audit & Separate research branch. & Complete action, inner product, ordering, power counting, time dependence and representation checks. No promised completion time.\\
GW geometry/tetrad explorer & Optional teaching extension. & Show coordinates versus physical distances honestly; geometry is not a mandatory precursor to the bar ODE.\\
Quantum oscillator and phonon probabilities & Conditional. & Classical benchmark passes; quantum question, occupation, cutoff, and preparation are specified.\\
Resonant spin-phonon swaps/avoided crossings & Conditional. & A nonzero matrix element, frequency match, coupling versus linewidths, and physical strain profile are demonstrated.\\
Dispersive phonon readout & Separate derivation. & A controlled detuned elimination and readout model; linear strain coupling alone is not QND energy measurement.\\
CPMG/XY8 & Add after one protocol. & Operator-specific filter, realistic pulses and consistent resource counting.\\
Multiple modes/FEM/anisotropic elasticity & Add when needed physically. & Analytic one-mode result and mesh/mode-truncation convergence first.\\
QFI and Bayesian inference & Conditional inference tools. & Specify observable likelihood, nuisance parameters and mixed-state dynamics as needed.\\
LIGO data/matched filtering & Optional data-analysis extension. & Separate measured detector noise from an incident waveform; source orientation and response are explicit.\\
WASM/Three.js and derivation map & Final presentation layer. & Native core has released benchmarks. Browser runs agree with native runs to documented tolerances.\\
ML, GPUs, autodifferentiation & Defer. & A measured bottleneck or a defined inference/optimization task exists.\\
Rarita--Schwinger extension & Remove from this plan. & It requires a distinct constrained higher-spin research project; it is not a few-day replacement $\bm S\mapsto\bm S^{(3/2)}$.\\
Single-graviton claims & Exclude from V1/V2 claims. & Quantized-source model and measurement argument that addresses classical alternatives; phonon counting alone is insufficient.\\
\bottomrule\end{longtable}

The proposed Rarita--Schwinger action has no specified mass term although its later expansion divides by $m$. Its canonical stress expression has uncontracted/inconsistently repeated indices, and ``spin terms'' is not an explicit Belinfante tensor. A massive vector-spinor requires constraints that remove unphysical components and a consistent background/coupling prescription. The cited Dirac work of Obukhov et al.~\cite{obukhov} does not supply that derivation. The original higher-spin and supergravity references are real~\cite{rarita,deser,supergravity}, but do not validate the formulas pasted here.

\section{Numerical and data rules that prevent false results}
\label{sec:numerics}

\begin{enumerate}
\item \textbf{Do not differentiate raw detector noise and call it a GW.} GWOSC strain is a calibrated detector-channel record containing noise, approximately $d=F_+\gamma_++F_\times\gamma_\times+n$ in the appropriate response limit. A single record does not determine both source polarizations. Use analytic signals first, then a documented model or posterior waveform. A noisy-channel injection is a separate signal-processing exercise.~\cite{gwosc,ligonoise}
\item \textbf{Differentiate deliberately.} For a smooth synthetic waveform use its analytic derivatives. A second finite difference amplifies sample noise roughly as $\Delta t^{-2}$ in amplitude; spectral differentiation multiplies by $-\omega^2$ and needs bandwidth and boundary control. Check sampling, monotonic times, units, interpolation, antialiasing and edge transients. Finite-difference convergence on a sine wave does not validate arbitrary noisy data.
\item \textbf{Use independent references.} Compare the ODE with Eq.~\eqref{eq:transfer} and convolution, not merely the same RK4 code translated into Python. At fixed smooth forcing, RK4 global error should decrease by about $16$ when the step is halved, until roundoff or interpolation error dominates. Repeat at resonances and other concrete risk points.
\item \textbf{A matrix exponential is exact only for constant $H$ on its stated interval.} Exponentiating a frozen time-dependent Hamiltonian is an integrator. A midpoint Cayley step
\[
 (I+i\Delta tH_{n+1/2}/2\hbar)\psi_{n+1}
 =(I-i\Delta tH_{n+1/2}/2\hbar)\psi_n
\]
is unitary for Hermitian $H_{n+1/2}$ but still has phase/time-discretization error. Norm preservation is not accuracy. Work in a documented rotating/interacting frame instead of resolving GHz carrier phases for a Hz-scale signal unnecessarily.
\item \textbf{Open-system checks include positivity.} For
$\dot\rho=-i[H,\rho]/\hbar+\sum_a(L_a\rho L_a^\dagger-\{L_a^\dagger L_a,\rho\}/2)$,
check trace, Hermiticity, eigenvalue nonnegativity within tolerance, and physical decay/thermal steady states. Generic RK4 need not preserve positivity. For $L=\sqrt\lambda S_z$, coherence $\rho_{mn}$ decays at $\lambda(m-n)^2/2$; one arbitrary dephasing coefficient is not the same $T_2$ for all qutrit coherences.
\item \textbf{Report error in the observable.} Population, phase, recovered signal amplitude, and inference error matter more than a solver's default tolerance. For tiny signals use rescaling or linear-response/variational equations; subtracting nearly identical double-precision populations can erase the effect. Test derivative step-size convergence for QFI; $10^{-21}$ finite differences are generally unusable without a parameter rescaling.
\item \textbf{Make PSD and FFT conventions explicit.} An FFT amplitude is not a PSD or an SNR. Record one/two-sided normalization, sample rate, window, equivalent noise bandwidth and frequency resolution. Avoid $q(f)/\gamma(f)$ where the input has negligible power; estimate transfer functions through controlled injections or an appropriate spectral estimator.
\item \textbf{Preserve provenance.} Every benchmark records units, configuration, source/DOI and equation, assumptions, software version, compiler/build, seed if applicable, and tested error tolerance. Distinguish established model, externally calibrated parameters, illustrative assumptions, and unverified derivations. Never label a formula ``verified'' merely because it runs.
\end{enumerate}

The frontend, if built, must use the simulated state, expose physical units, and label any visual displacement magnification or artificially amplified signal. A qutrit is not described by an ordinary three-component qubit Bloch vector. Population bars and selected coherences are a clearer first display. A ring of free masses can have fixed TT coordinates while proper separations change; the coordinate toggle must transform the quantities it shows.

\section{Audit register for the four supplied inputs}
\label{sec:register}

Input A is \emph{Innlimt tekst(20260918-093730).txt} (including its duplicated strategic-assessment/Ward sections); B is \emph{Innlimt tekst (2)(6).txt}; C is \emph{Innlimt tekst (3).txt}; D is the pasted \emph{Effective Spin-1 Gravitational Couplings for NV Centres} manuscript. This register complements the explicit derivations above; it does not treat repeated claims as independent evidence.

\begin{longtable}{@{}P{10mm}P{54mm}P{89mm}@{}}\toprule
Input & Claim or calculation & Finding and disposition\\\midrule\endhead
A--C & Tidal curvature and single-mode detector chain & Retain after restoring SI factors, mode normalization, finite-time response and a strain-transfer map. Sections~\ref{sec:conv}--\ref{sec:strain}.\\
A & $M=1800$ kg, $f_0=150$ Hz, $Q=10^{10}$ with unspecified $L$ & Not a complete design. Close the geometry/material relation; high $Q$ requires a provenance and finite-time calculation.\\
A & NV Rabi matrix in $\ket0,\ket{+1}$ & Incorrect: the stated quadrupole restricts to zero. Section~\ref{sec:nv}.\\
A & $10^{-11}$ Hz multiplied by $\sqrt N$ as Rabi enhancement & Product arithmetic is right; units and collective interpretation are wrong for independent classically driven spins.\\
A & $2p$ matrix element $3a_0^2/2$ and quoted splitting & Incorrect. $-12a^2$ in the standard phase convention, giving eigenvalues $0,\pm3\mu a^2\omega^2\gamma_+$.\\
A & $1s\to2p$ quadrupole absorption & Forbidden. Replace with an allowed $1s\to3d$ benchmark.\\
A & Gaunt formula and simple selection rules & Correct harmonic decomposition; incorrect conjugation in Gaunt formula. Add parity and triangle conditions and specify axis.\\
A & Golden Rule and $0.31\ell_{\rm Pl}^2$ & Missing cosine Fourier factor unless amplitude is redefined; cross section is a specific normalized frequency average, not a universal peak.\\
A & Spatial-only Ward proof and $p_ip_j$ identity & Both invalid. Replace with complete conservation and the force-dependent double commutator.\\
A & Spin part of EMT and automatic gauge cancellation & Invalid free-index structure. No conclusion follows from it.\\
A & Spin-$3/2$ action, EMT and substituted spin coupling & Internally incomplete/malformed and unsupported. Remove from execution scope.\\
A & Minor fixes make paper submittable in half a day & Unsupported. Novelty and submission readiness need review of the actual complete paper.\\
B & $\gamma=\kappa h^{\rm can}$ and separate model status & Retain. Add dimensions, empirical calibration status and conditions for verification.\\
B & Four FW ledger terms imply full expansion & Useful subset, not general completeness through $1/m^3$. Include time transformations and power counting.\\
B & Exact $4\times4$ Dirac comparison and gradient study & Conditional: homogeneous momentum block versus spatial/momentum-resolved propagating wave are different numerical problems.\\
B & $30a_0^2$, $-12a_0^2$, $\pm C$ hydrogen checks & Correct with the defined phases and $C$. Instantaneous eigenvalues are not a detection prediction.\\
B & Full geometry-to-sensor lab before first release & Excessive dependency chain. Retain as optional modules after mechanical V1.\\
B & Gauge interpolation slider & Teaching device only unless generated by an explicit transformation of Hamiltonian, state and observables.\\
C & Elastic wave equation & Correct for uniform one-dimensional linear elasticity without the drive; boundary conditions and material assumptions must be specified.\\
C & Zero-point normalization and $g_{\rm GW}$ & Mixed total/modal masses cause a $\sqrt2$ ambiguity; displayed drive coefficient has energy units. Distinguish GW and mode frequencies.\\
C & Mandatory zero-point-strain step in classical transduction & Conceptually wrong. Classical strain transfer exists without quantizing the mode.\\
C & Generic hybrid interaction drives every transition & Only with the required operator matrix elements, resonance/dressing and state preparation.\\
C & Trace and Hermiticity suffice for Lindblad validation & Add positivity, steady-state, decay-rate and truncation checks.\\
C & Pure-state QFI and $1/\sqrt{NF_Q}$ & Formula correct for pure probes and independent counted resources. No extra $N$ if already in the full state; not a mixed-state formula.\\
C & Twelve-week full programme; prewritten past-tense CV claims & Unjustified schedule. Use gates; describe work as completed only after it is implemented and validated.\\
D & Triplet basis and $^3A_2$ ground configuration & Retain for negatively charged NV$^-$ as a low-energy description. Triplet/singlet placement is not a substitute for the full many-body material model.\\
D & Vector projection intermediate value & $\langle\bm s_a\cdot\bm S\rangle=1$, not $1/2$; final $P\bm s_aP=\bm S/2$ is correct.\\
D & Joint-spin tensor projection & Correct. It does not derive the projection of an orbital kinetic tensor into spin space.\\
D & $\Xi\gamma_{ij}\{S_i,S_j\}$ to plus polarization & Loses a factor of two under the stated sum. Material coefficient remains unestablished.\\
D & Breaking TT gauge adds independent observable vector drive & Unsupported. Transform the full system and measurement; Fermi metric terms do not prove an NV coupling.\\
D & Trace mass modulation and gravitoelectric labels & Incorrect or insufficiently specified. See Section~\ref{sec:projection}.\\
D & $\Delta m_s$ uniquely fingerprints vector/tensor gravity & False. Rank-two spin operators can also have $\Delta m_s=\pm1$.\\
D & Universal microscopic-to-bar substitution & Unsupported shortcut. Derive continuum overlaps and normalized mode coordinates explicitly.\\
D & Dominant observable terms and rigorous foundation already established & Withdraw. No quantitative comparison, material matching or noise calculation establishes those claims.\\
\bottomrule\end{longtable}

\section{Source audit: identity and actual support}
\label{sec:sources}

The following checks distinguish a genuine reference from a reference that actually supports the attached claim. Bibliographic details and public full text were checked where relevant to retained equations; a real background citation is not a substitute for a missing derivation. This audit is dated 18 September 2026 and is not an exhaustive novelty search.

\begin{longtable}{@{}P{42mm}P{113mm}@{}}\toprule
Supplied source & Verified identity, relevance and correction\\\midrule\endhead
Abbott et al., PRL 116, 061102 (2016) & Correct: observation of GW150914.~\cite{abbott} Supports the historical detection, not an NV detector sensitivity.\\
Cai, Visinelli, Yan, PDU 53, 102374 (2026) & Real: \emph{Atomic quantum sensors for high-frequency gravitational wave searches}, DOI 10.1016/j.dark.2026.102374, arXiv:2510.15031. The checked v2 describes cavity electromagnetic transduction and atomic readout. It is not a derivation of direct NV spin-strain/gravitomagnetic coupling. Do not transfer its sensitivity forecasts to this project.~\cite{cai}\\
Doherty et al., Physics Reports 528, 1 (2013) & Correct; full page range 1--45.~\cite{doherty} Supports NV electronic structure and effective spin physics, not the claimed orbital-to-spin gravitational projection.\\
Fondo, thesis (2026) & The full thesis and its full nine-term derivation are not supplied in the four audited inputs. No complete title, repository record or independently checked derivation is supplied by this citation. Keep it as an internal dependency until supplied and audited; do not treat self-citation as verification.\\
Mashhoon, PRD 75, 084014 (2007) & Wrong attribution. That article concerns the running of Newton's constant and is by H. W. Hamber and R. M. Williams.~\cite{hamber} Replace the Fermi-coordinate support with Manasse and Misner~\cite{manasse}; a different Mashhoon paper must be identified if intended.\\
Tobar et al., arXiv:2401.xxxxx & Invalid placeholder and incorrect initials. Correct paper: G. Tobar, S. K. Manikandan, T. Beitel and I. Pikovski, Nature Communications 15, 7229 (2024), arXiv:2308.15440.~\cite{tobar} It proposes quantum acoustic detection; it is not an NV implementation or an achieved detector.\\
Ovartchaiyapong et al., Nat. Commun. 5, 4429 (2014) & Correct.~\cite{ovart} Full-text coefficient and sensitivity checks support the historical numbers used here, with their experimental conventions and uncertainty.\\
Teissier et al., PRL 113, 020503 (2014) & Correct: strain coupling to a diamond mechanical oscillator.~\cite{teissier} Supports mechanical NV coupling, not GW transduction.\\
MacQuarrie et al., PRL 111, 227602 (2013) & Correct: mechanically driven NV spin transitions.~\cite{macquarrie} Does not by itself specify the current resonator, local strain or incident-GW sensitivity.\\
Udvarhelyi et al., PRX 10, 011043 (2020) & Wrong journal/article attribution. That PRX article is \emph{Long-Range Prethermal Phases of Nonequilibrium Matter}. The relevant NV paper is PRB 98, 075201 (2018), arXiv:1712.02684.~\cite{udvarhelyi} Its full symmetry model contains six coefficients, including single-quantum strain terms.\\
Barfuss et al., Nature Physics 11, 820 (2015) & Correct: \emph{Strong mechanical driving of a single electron spin}, pages 820--824.~\cite{barfuss} Supports strain-driven control, not the proposed gravitational derivation.\\
Taylor et al., Nature Physics 4, 810 (2008) & Correct: \emph{High-sensitivity diamond magnetometer with nanoscale resolution}, pages 810--816.~\cite{taylor} Magnetometry sensitivity is not automatically strain sensitivity.\\
S. A. Zaidi et al., arXiv:2510.11009 & Wrong authors. It is J. Liang, M. Li, Y. Gao, W. Ji, S. Sun and Q.-S. Yan, \emph{Detecting gravitational waves with spin systems} (2025).~\cite{liang} It discusses a distinct spin/NMR proposal; its forecast does not validate an NV-mechanical detector.\\
Boughn and Rothman, CQG 23, 5839 (2006) & Correct, pages 5839--5852.~\cite{boughn} Their exact cross-section definition and allowed transition matter; the input's $1s\to2p$ sketch does not reproduce their result.\\
Quach (2015); Goncalves et al. (2007) & Quach arXiv:1509.01671 explicitly corrects an EM gauge-variant FW result. The earlier paper is B. Goncalves, Y. N. Obukhov and I. L. Shapiro, PRD 75, 124023 (2007).~\cite{quach,goncalves} The submitted Ward proof is neither a derivation nor a faithful summary of that correction.\\
Rarita--Schwinger (1941), Deser--Zumino (1976), Freedman--Van Proeyen (2012) & Real foundational higher-spin/supergravity sources.~\cite{rarita,deser,supergravity} Their identities do not validate the malformed action/EMT or massive FW analogue in input A.\\
Obukhov--Silenko--Teryaev, PRD 80, 064044 (2009) & Correct title concerns spin dynamics in fields of rotating bodies and the equivalence principle.~\cite{obukhov} It studies Dirac/classical spin and tetrad dependence, not the asserted Rarita--Schwinger reduction.\\
\bottomrule\end{longtable}

\section{Deliverables and claims allowed at each completion point}

\begin{tabularx}{\textwidth}{@{}P{24mm}XX@{}}\toprule
Release & Concrete deliverable & Claim supported\\\midrule
V1 & Native C++ bar-response executable, analytic checks, finite-signal results and documented noise assumptions. & ``Validated a reduced mechanical GW-response model.''\\
V2 & Strain-calibrated NV model, readout likelihood and a consistent transfer/noise study. & ``Evaluated conditional NV readout feasibility for the specified structure.''\\
Research note & Full corrected derivation, source comparison, representation checks and uncertainty in novelty. & Only the results actually proved and reproduced.\\
Frontend & Native/WASM agreement and clearly labelled numerical visualizations. & An interactive view of the validated model.\\
\bottomrule\end{tabularx}

The original claims of a new rigorous spin-1 gravitational Hamiltonian, experimentally relevant dominant terms, automatic collective Rabi enhancement, and a ready-to-submit paper are not supported by the supplied material. The corrected programme can be completed in useful stages without assuming any of those claims.

\appendix
\section{Reproducible calculation checks supplied with this plan}

\code{verify\_gw\_plan.py} checks the finite spin algebra, triplet projections, hydrogen radial/angular results, bar normalization and overlap, the oscillator counterexample to the false momentum identity, numerical scales, the bar transfer function against independent time evolution, RK4 convergence, and commuting/noncommuting QFI benchmarks. Results are recorded in \code{calculation\_results.json}. The checks are specific to the reduced models; they do not certify a complete FW calculation, an experimental device, a paper's novelty, or the full correctness of every cited paper.

The accompanying run passed all 41 scripted assertions. The independent harmonic-response comparison had maximum absolute displacement error $4.98\times10^{-12}$ in its dimensionless test; RK4 step-halving error ratios were $16.012$ and $16.007$. Exact symbolic checks were used for the spin projections, hydrogen integrals and mode normalization. These results validate the listed benchmark calculations, not the missing material or relativistic derivations.

Run the script with the Python dependencies stated in its header. The LaTeX file contains its own bibliography and can be compiled with \code{pdflatex} twice. No external figure files are required. Numeric inputs labelled illustrative must remain labelled illustrative when reused.

\begin{thebibliography}{99}\small
\bibitem{abbott} B. P. Abbott et al., \emph{Observation of Gravitational Waves from a Binary Black Hole Merger}, Phys. Rev. Lett. \textbf{116}, 061102 (2016). \href{https://doi.org/10.1103/PhysRevLett.116.061102}{doi:10.1103/PhysRevLett.116.061102}.
\bibitem{manasse} F. K. Manasse and C. W. Misner, \emph{Fermi Normal Coordinates and Some Basic Concepts in Differential Geometry}, J. Math. Phys. \textbf{4}, 735--745 (1963). \href{https://doi.org/10.1063/1.1724316}{doi:10.1063/1.1724316}.
\bibitem{tobar} G. Tobar, S. K. Manikandan, T. Beitel and I. Pikovski, \emph{Detecting single gravitons with quantum sensing}, Nat. Commun. \textbf{15}, 7229 (2024). \href{https://doi.org/10.1038/s41467-024-51420-8}{doi:10.1038/s41467-024-51420-8}; \href{https://arxiv.org/abs/2308.15440}{arXiv:2308.15440}.
\bibitem{doherty} M. W. Doherty et al., \emph{The nitrogen-vacancy colour centre in diamond}, Phys. Rep. \textbf{528}, 1--45 (2013). \href{https://doi.org/10.1016/j.physrep.2013.02.001}{doi:10.1016/j.physrep.2013.02.001}.
\bibitem{udvarhelyi} P. Udvarhelyi, V. O. Shkolnikov, A. Gali, G. Burkard and A. Palyi, \emph{Spin-strain interaction in nitrogen-vacancy centers in diamond}, Phys. Rev. B \textbf{98}, 075201 (2018). \href{https://doi.org/10.1103/PhysRevB.98.075201}{doi:10.1103/PhysRevB.98.075201}; \href{https://arxiv.org/abs/1712.02684}{arXiv:1712.02684}, Eq. (3), Table I.
\bibitem{ovart} P. Ovartchaiyapong, K. W. Lee, B. A. Myers and A. C. Bleszynski Jayich, \emph{Dynamic strain-mediated coupling of a single diamond spin to a mechanical resonator}, Nat. Commun. \textbf{5}, 4429 (2014). \href{https://doi.org/10.1038/ncomms5429}{doi:10.1038/ncomms5429}; \href{https://arxiv.org/abs/1403.4173}{arXiv:1403.4173}.
\bibitem{teissier} J. Teissier, A. Barfuss, P. Appel, E. Neu and P. Maletinsky, \emph{Strain Coupling of a Nitrogen-Vacancy Center Spin to a Diamond Mechanical Oscillator}, Phys. Rev. Lett. \textbf{113}, 020503 (2014). \href{https://doi.org/10.1103/PhysRevLett.113.020503}{doi:10.1103/PhysRevLett.113.020503}.
\bibitem{macquarrie} E. R. MacQuarrie et al., \emph{Mechanical Spin Control of Nitrogen-Vacancy Centers in Diamond}, Phys. Rev. Lett. \textbf{111}, 227602 (2013). \href{https://doi.org/10.1103/PhysRevLett.111.227602}{doi:10.1103/PhysRevLett.111.227602}.
\bibitem{barfuss} A. Barfuss, J. Teissier, E. Neu, A. Nunnenkamp and P. Maletinsky, \emph{Strong mechanical driving of a single electron spin}, Nat. Phys. \textbf{11}, 820--824 (2015). \href{https://doi.org/10.1038/nphys3411}{doi:10.1038/nphys3411}.
\bibitem{taylor} J. M. Taylor et al., \emph{High-sensitivity diamond magnetometer with nanoscale resolution}, Nat. Phys. \textbf{4}, 810--816 (2008). \href{https://doi.org/10.1038/nphys1075}{doi:10.1038/nphys1075}.
\bibitem{boughn} S. Boughn and T. Rothman, \emph{Aspects of graviton detection: graviton emission and absorption by atomic hydrogen}, Class. Quantum Grav. \textbf{23}, 5839--5852 (2006). \href{https://doi.org/10.1088/0264-9381/23/20/006}{doi:10.1088/0264-9381/23/20/006}; \href{https://arxiv.org/abs/gr-qc/0605052}{arXiv:gr-qc/0605052}.
\bibitem{quach} J. Q. Quach, \emph{Foldy-Wouthuysen transformation of the generalised Dirac Hamiltonian in a gravitational-wave background} (2015). \href{https://arxiv.org/abs/1509.01671}{arXiv:1509.01671}, especially Eqs. (8)--(9), (22).
\bibitem{goncalves} B. Goncalves, Y. N. Obukhov and I. L. Shapiro, \emph{Exact Foldy-Wouthuysen transformation for gravitational waves and magnetic field background}, Phys. Rev. D \textbf{75}, 124023 (2007). \href{https://doi.org/10.1103/PhysRevD.75.124023}{doi:10.1103/PhysRevD.75.124023}.
\bibitem{obukhov} Y. N. Obukhov, A. J. Silenko and O. V. Teryaev, \emph{Spin dynamics in gravitational fields of rotating bodies and the equivalence principle}, Phys. Rev. D \textbf{80}, 064044 (2009). \href{https://doi.org/10.1103/PhysRevD.80.064044}{doi:10.1103/PhysRevD.80.064044}.
\bibitem{cai} Y.-F. Cai, L. Visinelli and S.-F. Yan, \emph{Atomic quantum sensors for high-frequency gravitational wave searches}, Phys. Dark Univ. \textbf{53}, 102374 (2026). \href{https://doi.org/10.1016/j.dark.2026.102374}{doi:10.1016/j.dark.2026.102374}; \href{https://arxiv.org/abs/2510.15031}{arXiv:2510.15031}, v2 checked.
\bibitem{liang} J. Liang, M. Li, Y. Gao, W. Ji, S. Sun and Q.-S. Yan, \emph{Detecting gravitational waves with spin systems} (2025). \href{https://arxiv.org/abs/2510.11009}{arXiv:2510.11009}.
\bibitem{carney} D. Carney, V. Domcke and N. L. Rodd, \emph{Graviton detection and the quantization of gravity}, Phys. Rev. D \textbf{109}, 044009 (2024). \href{https://doi.org/10.1103/PhysRevD.109.044009}{doi:10.1103/PhysRevD.109.044009}; \href{https://arxiv.org/abs/2308.12988}{arXiv:2308.12988}.
\bibitem{braunstein} S. L. Braunstein and C. M. Caves, \emph{Statistical distance and the geometry of quantum states}, Phys. Rev. Lett. \textbf{72}, 3439--3443 (1994). \href{https://doi.org/10.1103/PhysRevLett.72.3439}{doi:10.1103/PhysRevLett.72.3439}.
\bibitem{degen} C. L. Degen, F. Reinhard and P. Cappellaro, \emph{Quantum sensing}, Rev. Mod. Phys. \textbf{89}, 035002 (2017). \href{https://doi.org/10.1103/RevModPhys.89.035002}{doi:10.1103/RevModPhys.89.035002}.
\bibitem{gwosc} Gravitational Wave Open Science Center, \emph{Tutorials and guide to GW detections and noise}. \href{https://gwosc.org/tutorials/}{gwosc.org/tutorials}, accessed 18 September 2026.
\bibitem{ligonoise} B. P. Abbott et al., \emph{A guide to LIGO-Virgo detector noise and extraction of transient gravitational-wave signals}, Class. Quantum Grav. \textbf{37}, 055002 (2020). \href{https://arxiv.org/abs/1908.11170}{arXiv:1908.11170}.
\bibitem{hamber} H. W. Hamber and R. M. Williams, \emph{Renormalization group running of Newton's constant G: The static isotropic case}, Phys. Rev. D \textbf{75}, 084014 (2007). \href{https://doi.org/10.1103/PhysRevD.75.084014}{doi:10.1103/PhysRevD.75.084014}. Bibliographic correction only.
\bibitem{rarita} W. Rarita and J. Schwinger, \emph{On a Theory of Particles with Half-Integral Spin}, Phys. Rev. \textbf{60}, 61 (1941). \href{https://doi.org/10.1103/PhysRev.60.61}{doi:10.1103/PhysRev.60.61}. Background only.
\bibitem{deser} S. Deser and B. Zumino, \emph{Consistent supergravity}, Phys. Lett. B \textbf{62}, 335--337 (1976). \href{https://doi.org/10.1016/0370-2693(76)90089-7}{doi:10.1016/0370-2693(76)90089-7}. Background only.
\bibitem{supergravity} D. Z. Freedman and A. Van Proeyen, \emph{Supergravity}, Cambridge University Press (2012). \href{https://doi.org/10.1017/CBO9781139026833}{doi:10.1017/CBO9781139026833}. Background only.
\end{thebibliography}
\end{document}
